\newif\ifglobal

%%%%%%%%%%%%%%%%%%%%%%%
%%% Compile to view %%%
%%%%%%%%%%%%%%%%%%%%%%%

%%%%%%%%%%%%%%%%%%%%%%%%%%%%%%%%%%%%%%%%%%%%%%%%%%%
%%% Readme:                                     %%%
%%% https://www.overleaf.com/read/bwdjvfjksvjy/ %%%
%%%%%%%%%%%%%%%%%%%%%%%%%%%%%%%%%%%%%%%%%%%%%%%%%%%

% >>> M?: marks a module setting number or important notice

% >>> M4: Set {./relative-path} to external files for this module <<<
\def \modulefiles {./10-SYSTIMER}
% To add an external file, use:
% (Replace '---' with actual filename)
% '\input{\modulefiles/tables/---.tex}' for tables
% '\includegraphics{\modulefiles/figures/---}' for figures
% '\subfile{\modulefiles/---}' for register files


% >>> M1: This module is in global project? <<<
% 'YES' - uncomment; 'NO' - comment out
\globaltrue

\ifglobal
    \documentclass[main\_\_CN.tex]{subfiles}

\else
    % Font "FandolSong-Regular" used by ctex does not have GSUB table.
    % It causes warnings that do not affect compilation and can be ignored.
    \PassOptionsToPackage{quiet}{fontspec}

    \documentclass[openany, 10pt]{book}

    % >>> M2: In ./00-shared/preamble-trm-module, also find
    %         settings for visibility of todo notes, watermark <<<
    \usepackage{preamble-trm-module}


    % Enable Chinese typesetting
    \usepackage{ctex}

    % >>> M3: Temporary labels in Chinese
    %         you can add your own labels there <<<
    \myexternaldocument{./00-shared/temp-labels-cn}

\fi %\ifglobal


\setcounter{secnumdepth}{4}


% >>> M5: If this module needs any updates in
% './00-shared/preamble-trm-module.sty'
% (include additional package, etc.), contact your module PM
% or create the file './00-shared/changelog.tex' make the
% necessary updates yourself and document them in this file <<<

% Make text left-aligned and allow latex to break pages naturally, comment out for full justification
\raggedright
\raggedbottom

\begin{document}


% Sets language into which pre-defined bits of text are translated
\selectlanguage{chinese}
\setchineseformatting

% Do not delete this block
\ifglobal
% Add the following front matter if
% this module is outside of global project
\else
    \InputIfFileExists{readme}{}{}
    \listoftodos
    {\let\clearpage\relax \listofdonetodos}
    \tableofcontents
    %\listoftables
    %\listoffigures
\fi



%%%%%%%%%%%%%%%%%%%%%%%%%
%%% TRM Module Begins %%%
%%%%%%%%%%%%%%%%%%%%%%%%%

\hypertarget{systimer}{}
\chapter{系统定时器 (SYSTIMER)}
\label{mod:systimer}


\section{概述}

ESP32-S2 芯片内置一个 64 位系统定时器。该系统定时器专门用于操作系统，可生成周期性系统滴答或延时中断用于调度系统任务。

\section{主要特征}
\begin{itemize}
    \item 采用 64 位定时器。
    \item 使用 APB\_CLK 时钟。
    \item 每个 APB\_CLK 周期，定时器数值的递增幅度可配置。
    \item 如果 APB\_CLK 时钟源发生切换，包括从 PLL\_CLK 切换至 XTAL\_CLK，或从 XTAL\_CLK 切换至 PLL\_CLK，系统定时器将执行自动时间补偿，以提高定时器准确性。
    \item 支持设置不同的报警值或周期（目标）生成三个独立中断。
    \item 支持 64 位报警值和 30 位报警周期。
    \item 支持软件配置基准计数值。例如，支持从 Light-sleep 唤醒之后，系统定时器通过软件加载 RTC 定时器记录的睡眠时间，并进行补偿。
    \item CPU 处于停止状态或处于在线调试状态时，系统定时器也将停止运行。
\end{itemize}

\section{时钟源选择}

用户可从 XTAL\_CLK 或 PLL\_CLK 中选择一个作为时钟源。选择方式见章节 \ref{mod:resclk} \textit{\nameref{mod:resclk}}中的表 \ref{table:cpu clksrc} \textit{\nameref{table:cpu clksrc}}。具体用于系统定时器的时钟频率见表 \ref{table:apb_clk_src} \textit{\nameref{table:apb_clk_src}}。定时器将根据使用的具体时钟源，每一时钟周期按照 \hyperref[fielddesc:SYSTIMERTIMERXTALSTEP]{SYSTIMER\_TIMER\_XTAL\_STEP} 或 \hyperref[fielddesc:SYSTIMERTIMERPLLSTEP]{SYSTIMER\_TIMER\_PLL\_STEP} 中设定的幅度递增。


\section{功能描述}

下图 \ref{figure:System Timer Structure} 为系统定时器结构图。用户可通过置位寄存器 \hyperref[regdesc:SYSTEMPERIPCLKEN0REG]{SYSTEM\_PERIP\_CLK\_EN0\_REG} 中的 \hyperref[regdesc:SYSTEMPERIPCLKEN0REG]{SYSTEM\_SYSTIMER\_CLK\_EN} 位使能系统定时器，或置位寄存器 \hyperref[regdesc:SYSTEMPERIPRSTEN0REG]{SYSTEM\_PERIP\_RST\_EN0\_REG} 中的 \hyperref[regdesc:SYSTEMPERIPRSTEN0REG]{SYSTEM\_SYSTIMER\_RST} 位通过软件复位系统定时器。详细信息可参考章节 \ref{mod:sysreg} \textit{\nameref{mod:sysreg}}中表 \ref{table:peripheralcontrollingbit} \textit{\nameref{table:peripheralcontrollingbit}}。

\begin{figure}[H]
    \centering
    \includegraphics[width=0.9\textwidth]{\modulefiles/figures/system_timer.eps}
    \caption{系统定时器结构}
    \label{figure:System Timer Structure}
\end{figure}

\subsection{读取系统定时器的值} \label{section:read-timer-value}
\begin{enumerate}
    \item 置位 \hyperref[fielddesc:SYSTIMERTIMERUPDATE]{SYSTIMER\_TIMER\_UPDATE} 将系统定时器的值更新至寄存器。
    \item 等待 \hyperref[fielddesc:SYSTIMERTIMERVALUEVALID]{SYSTIMER\_TIMER\_VALUE\_VALID} 置位。置位后则可以从寄存器中读取系统定时器的值。
    \item 从 \hyperref[fielddesc:SYSTIMERTIMERVALUEHI]{ SYSTIMER\_TIMER\_VALUE\_HI} 中读取定时器数值的高 32 位，从 \hyperref[fielddesc:SYSTIMERTIMERVALUELO]{SYSTIMER\_TIMER\_VALUE\_LO} 中读取定时器数值的低 32 位。
\end{enumerate}

\subsection{设置一次性延时报警}
\begin{enumerate}
    \item 读取当前系统定时器的值，步骤见章节 \ref{section:read-timer-value}。读取的当前值可用于计算步骤 3 中的报警值。
    \item 清零 \hyperref[fielddesc:SYSTIMERTARGET0PERIODMODE]{ SYSTIMER\_TARGET\regindex{x}\_PERIOD\_MODE}，设置定时器工作模式为一次性延时报警。
    \item 将目标值（报警值）高 32 位写入 \hyperref[fielddesc:SYSTIMERTIMERTARGET0HI]{SYSTIMER\_TIMER\_TARGET\regindex{x}\_HI}，低 32 位写入 \hyperref[fielddesc:SYSTIMERTIMERTARGET0LO]{SYSTIMER\_TIMER\_TARGET\regindex{x}\_LO}。
    \item 置位 \hyperref[fielddesc:SYSTIMERTARGET0WORKEN]{ SYSTIMER\_TARGET\regindex{x}\_WORK\_EN}，使能上述选择的工作模式。
    \item 置位 \hyperref[fielddesc:SYSTIMERINT0ENA]{SYSTIMER\_INT\regindex{x}\_ENA} 使能定时器中断。定时器计数达到报警值则触发中断。
\end{enumerate}

\subsection{设置周期性报警}
\begin{enumerate}
    \item 置位 \hyperref[fielddesc:SYSTIMERTARGET0PERIODMODE]{ SYSTIMER\_TARGET\regindex{x}\_PERIOD\_MODE}，设置定时器工作模式为周期性报警。
    \item 将目标值（报警周期）写入 \hyperref[fielddesc:SYSTIMERTARGET0PERIOD]{SYSTIMER\_TARGET\regindex{x}\_PERIOD}。
    \item 置位 \hyperref[fielddesc:SYSTIMERTARGET0WORKEN]{ SYSTIMER\_TARGET\regindex{x}\_WORK\_EN}，使能上述选择的工作模式。
    \item 置位 \hyperref[fielddesc:SYSTIMERINT0ENA]{SYSTIMER\_INT\regindex{x}\_ENA} 使能定时器中断。定时器计数达到目标值则触发中断。
\end{enumerate}

\subsection{定时器更新与睡眠模式}
\begin{enumerate}
    \item 在芯片进入 Light-sleep 模式之前，用户需配置 RTC 定时器用于精确记录睡眠时间，见低功耗管理章节。
    \item 系统从 Light-sleep 模式唤醒后开始读取 RTC 定时器记录的睡眠时间，见低功耗管理章节。
    \item 读取系统定时器当前值，具体步骤见章节 \ref{section:read-timer-value}。
    \item 定时器当前值与睡眠时间相加。
    \item 相加的和写入 \hyperref[fielddesc:SYSTIMERTIMERLOADHI]{SYSTIMER\_TIMER\_LOAD\_HI}（高 32 位）和 \hyperref[fielddesc:SYSTIMERTIMERLOADLO]{SYSTIMER\_TIMER\_LOAD\_LO}（低 32 位）。
    \item 置位 \hyperref[fielddesc:SYSTIMERTIMERLOAD]{SYSTIMER\_TIMER\_LOAD} 将步骤 5 中寄存器的值加载至系统定时器，实现系统定时器的更新。
\end{enumerate}

\section{基地址}\label{system-timer-base-address}
用户可以通过两个不同的寄存器基地址访问系统定时器相关的寄存器，如表 \ref{table:systimer_base_address} 所示。更多信息，请访问章节 \ref{mod:sysmem} \textit{\nameref{mod:sysmem}}。


\begin{longtable}[c]{|l|C{2.5cm}| }
        \caption{系统定时器基地址}
		\label{table:systimer_base_address} \\
	    \hline
        \rowcolor{lightgray}
        访问方式 & 基地址 \\
        \hline
        PeriBUS1 & \hexprefix{}3F423000 \\
        \hline
        PeriBUS2 & \hexprefix{}60023000 \\
        \hline
\end{longtable}

\hypertarget{systimer-reg-summ}{}
\section{寄存器列表}\label{register summary-30}

请注意，表中的地址为相对于基地址的地址偏移量（相对地址）。请参考章节 \ref{system-timer-base-address} 获取基地址相关信息。

\subfile{\modulefiles/system_timer_summary_cn.tex}


\section{寄存器}\label{registers-30}
\subfile{\modulefiles/system_timer_details_cn.tex}

\end{document}

